\chapter{KẾT QUẢ THỰC TẬP}

\section{Tóm tắt nội dung đã thực hiện}

Sinh viên tóm tắt toàn bộ các nội dung đã thực hiện trong quá trình thực tập, bao gồm các công việc chuyên môn, công việc hỗ trợ, các buổi họp, trao đổi nghiệp vụ hoặc các hoạt động khác tại doanh nghiệp.

\section{Sản phẩm hoặc dự án thực tế đã tham gia}

Sinh viên mô tả sản phẩm, module, chức năng hoặc dự án thực tế đã tham gia. Nội dung có thể bao gồm:

\begin{itemize}
    \item Tên dự án hoặc sản phẩm.
    \item Mục tiêu của dự án.
    \item Vai trò của sinh viên trong dự án.
    \item Các chức năng hoặc nhiệm vụ đã thực hiện.
    \item Công nghệ, ngôn ngữ lập trình, framework, cơ sở dữ liệu hoặc công cụ sử dụng.
\end{itemize}

\begin{table}[H]
\centering
\caption{Danh sách công nghệ sử dụng trong quá trình thực tập}
\label{tab:cong-nghe-su-dung}
\begin{tabular}{|c|p{4cm}|p{8cm}|}
\hline
\textbf{STT} & \textbf{Công nghệ/Công cụ} & \textbf{Mục đích sử dụng} \\
\hline
1 & Git/GitHub & Quản lý mã nguồn và phối hợp làm việc nhóm. \\
\hline
2 & Visual Studio Code & Soạn thảo mã nguồn. \\
\hline
3 & MySQL/\allowbreak PostgreSQL & Lưu trữ và quản lý dữ liệu. \\
\hline
\end{tabular}
\Nguon{Sinh viên tổng hợp}
\end{table}

\section{Kỹ năng đã rèn luyện được}

\subsection{Kỹ năng chuyên môn}

\begin{itemize}
    \item Kỹ năng lập trình, kiểm thử hoặc phân tích hệ thống.
    \item Kỹ năng sử dụng công cụ phát triển phần mềm.
    \item Kỹ năng đọc hiểu tài liệu kỹ thuật.
    \item Kỹ năng xử lý lỗi và tối ưu công việc.
\end{itemize}

\subsection{Kỹ năng mềm}

\begin{itemize}
    \item Kỹ năng giao tiếp trong môi trường làm việc.
    \item Kỹ năng làm việc nhóm.
    \item Kỹ năng quản lý thời gian.
    \item Tác phong làm việc chuyên nghiệp.
\end{itemize}

\section{Đề xuất cải thiện hoặc đóng góp ý kiến cho doanh nghiệp}

Sinh viên nêu các đề xuất phù hợp nếu có, ví dụ:

\begin{itemize}
    \item Đề xuất cải thiện quy trình hướng dẫn thực tập sinh.
    \item Đề xuất cải thiện tài liệu kỹ thuật hoặc tài liệu bàn giao.
    \item Đề xuất cải thiện công cụ, quy trình kiểm thử hoặc quy trình triển khai.
\end{itemize}

\section{Tự đánh giá hiệu quả thực tập}

Sinh viên tự đánh giá mức độ hoàn thành công việc, những điểm mạnh, hạn chế và định hướng cải thiện sau quá trình thực tập.
